% SPDX-License-Identifier: Apache-2.0
\section{One Bundle for Both Its Ends}
\label{sec:x-flipped}

A bundle that runs both ways, a request one way and the answer the
other, was two structs: every \code{In} of one an \code{Out} of the
other, or a \code{Tx} and an \code{Rx}, written twice and kept in step
by hand. A struct generic over its side is written once (issue 1650).
A field \code{S::Drive<T>} is an \code{Out<T>} on the
\code{Driver}'s side and an \code{In<T>} on the \code{Receiver}'s,
\code{S::Read<T>} the reverse, and \code{S::Send<T>} and
\code{S::Recv<T>} are a channel's two ends the same way.
\code{\#[derive(Ports)]} writes \code{Link} for both sides, so
\code{link::<Dbg<Receiver>>()} makes the two ends at once, as
Section~\ref{sec:x-link} makes a link of two structs.

Here \code{Dm} takes \code{Dbg<Driver>} and raises \code{haltreq}
every other cycle, and \code{Hart} takes \code{Dbg<Receiver>}, sends
its count on \code{rdata} and reports the request it reads.
\code{Pair} joins them with one link. The answer is a channel rather
than a wire: \code{Dm} runs before \code{Hart} in a step, and a wire is
set before it is read, while a channel is registered. The run checks
both reports cycle by cycle and prints the ports of the two sides, the
same names each turned round; the netlist is checked against the run
under both simulators. AXI-Lite's \code{LitePort} and
\code{LiteHostPort} are now one such struct, \code{LiteBus}, with the
same netlists.

\begin{figure}[htbp]
\centering
\input{flipped_timing.tex}
\caption{The request every other cycle, the count, and what each side
reports of the other.}
\label{fig:x-flipped}
\end{figure}

\lstinputlisting[caption={\code{ex\_flipped.rs}. Run with \code{bazel run //lib/examples:ex\_flipped}.},label={lst:x-flipped}]{lib/examples/ex_flipped.rs}

\lstinputlisting[style=ascii,caption={What \code{ex\_flipped} printed when the build ran it: the two reports each cycle, the ports of each side, then the lowering.},label={lst:x-flipped-out}]{out_flipped.txt}
